%=====================================================================
\section{Inside the switch}
\label{sec:switch}

\subsection{Life of a packet through the pipeline}

\begin{figure}[H]\centering
\begin{tikzpicture}
  \node[term] (in) at (0,0) {frame arrives on port 1 or 2 (or from the controller, port 255)};
  \node[dec] (inst) at (0,-1.4) {second pass\\ (resubmit)?};
  \node[swbox, text width=34mm] (eth) at (5,-1.4) {parse Ethernet only:\\ every other byte stays payload};
  \node[dec] (cpu) at (0,-2.9) {from port 255?};
  \node[swbox, text width=34mm] (po) at (5,-2.9) {read \code{packet\_out}: reinjected = 1,\\ kind, certificate index};
  \node[swbox, text width=50mm] (l34) at (0,-4.3) {parse Ethernet (up to two VLAN tags), IPv4\\ (+options) or IPv6, TCP and every option byte, or UDP};
  \node[dec] (pay) at (0,-5.8) {message in\\ the payload?};
  \node[swbox, text width=34mm] (walk) at (-4.3,-7.3) {whole message: walk it\\ (TLS hello, HTTP request,\\ DHCP, certificate):\\ values $\to$ metadata};
  \node[swbox, text width=34mm] (inc) at (4.3,-7.3) {message starts here but\\ continues later:\\ mark \code{incomplete}};
  \node[swbox, text width=64mm] (ing) at (0,-9.0) {\textbf{Ingress}: forward table, ConnKey, ConnId, Handshake, Stream (copy?),\\ Tls, Http, Ssh, Dhcp, X509 $\to$ digests to the controller};
  \node[dec] (wk) at (0,-10.6) {bytes consumed, or\\ a 2nd digest due?};
  \node[swbox, text width=34mm] (rs) at (5,-10.6) {\code{resubmit}: the original\\ packet runs again, keeping\\ inPort, connection number\\ and the handshake record};
  \node[swbox, text width=64mm] (eg) at (0,-12.1) {\textbf{Egress}: if \code{copy}: clone (E2E) to session 100, the CPU port;\\ on the clone: fill \code{packet\_in} (offset, direction, start, type)};
  \node[swbox, text width=64mm] (dp) at (0,-13.4) {\textbf{Deparser}: emit the extracted headers, then the payload};
  \node[term] (o1) at (-2.8,-14.7) {out of the other port,\\ byte for byte unchanged};
  \node[term] (o2) at (2.8,-14.7) {copy to the controller\\ (packet-in)};
  \draw[arr] (in) -- (inst);
  \draw[arr] (inst) -- node[lbl, above]{yes} (eth);
  \draw[arr] (inst) -- node[lbl, right]{no} (cpu);
  \draw[arr] (cpu) -- node[lbl, above]{yes} (po);
  \draw[arr] (po.south) |- ($(l34.east)+(0,0)$);
  \draw[arr] (cpu) -- node[lbl, right]{no} (l34);
  \draw[arr] (l34) -- (pay);
  \draw[arr] (pay) -| node[lbl, above, pos=0.3]{whole} (walk);
  \draw[arr] (pay) -| node[lbl, above, pos=0.3]{partial} (inc);
  \draw[arr] (pay) -- node[lbl]{none} (ing);
  \draw[arr] (walk) |- ($(ing.north west)+(6mm,0)$) -- ++(0,-0.01);
  \draw[arr] (inc) |- ($(ing.north east)+(-6mm,0)$) -- ++(0,-0.01);
  \draw[arr] (eth.east) -- (7.4,-1.4) |- ($(ing.east)+(0,0)$);
  \draw[arr] (ing) -- (wk);
  \draw[arr] (wk) -- node[lbl, above]{yes} (rs);
  \draw[arr] (rs.east) -- (8.0,-10.6) |- (in.east);
  \draw[arr] (wk) -- node[lbl, right]{no} (eg);
  \draw[arr] (eg) -- (dp); \draw[arr] (dp) -- (o1); \draw[arr] (dp) -- (o2);
\end{tikzpicture}
\caption{One packet's path (\code{p4/parser.p4}, \code{p4/main.p4}). A walk reads a message with \code{advance},
which removes the read bytes from the packet: such a packet is \emph{resubmitted}, and on the second pass only
Ethernet is parsed, so the original bytes are forwarded. A packet from the controller (port 255) is read the same
way but never forwarded.}
\label{fig:life}
\end{figure}

\why{\code{advance} is the only way for a v1model parser to skip bytes it does not extract; extracting every byte
into headers would need header stacks thousands of entries long and the deparser would have to re-emit all of
them. Skipping and forwarding the untouched original on a second pass keeps the walk cheap and the forwarding
exact (checked frame by frame: 9,013 of 9,013).}
\whynot{\emph{Clone first, then walk the clone:} an ingress clone is parsed again from the start, so it would
lose its bytes the same way. \emph{Recirculate:} goes through egress and the deparser, which would emit the already
shortened packet.}

%---------------------------------------------------------------------
\subsection{The parser: top-level state machine}

\begin{figure}[H]\centering
\begin{tikzpicture}
  \node[st] (start) at (0,0) {start};
  \node[st] (ethonly) at (4,0) {parse\_ethernet\_only\\ (second pass)};
  \node[st, fill=white] (acc) at (8,0) {accept};
  \node[st] (po) at (0,-1.2) {parse\_packet\_out};
  \node[st] (eth) at (0,-2.4) {parse\_ethernet};
  \node[st] (vl) at (-5.8,-2.4) {parse\_vlan\_outer /\\ parse\_vlan\_inner};
  \node[st] (v4) at (-2.6,-3.6) {parse\_ipv4 (+options)};
  \node[st] (v6) at (2.6,-3.6) {parse\_ipv6};
  \node[st] (tcp) at (-2.6,-4.9) {parse\_tcp};
  \node[st] (udp) at (2.6,-4.9) {parse\_udp};
  \node[st] (opt) at (-2.6,-6.1) {TCP option walk (Fig.~\ref{fig:opt})};
  \node[st] (d4) at (1.4,-6.1) {DHCPv4 walk (67/68)};
  \node[st] (d6) at (4.6,-6.1) {DHCPv6 walk (546/547)};
  \node[st] (tp) at (-2.6,-7.3) {tcp\_payload\\ (``SSH-''? noted)};
  \node[st] (x) at (-6.4,-8.7) {x509\_start\\ (Fig.~\ref{fig:x509})};
  \node[st] (short) at (-2.6,-8.7) {tls\_short};
  \node[st] (peek) at (1.4,-8.7) {tls\_peek};
  \node[st] (inc) at (-2.6,-10.1) {tls\_incomplete /\\ http\_incomplete};
  \node[st] (tlschk) at (1.4,-10.1) {tls\_check $\to$ TLS walk};
  \node[st] (httpchk) at (5.4,-10.1) {http\_check $\to$ HTTP walk};
  \draw[arr] (start) -- node[lbl]{resubmit} (ethonly);
  \draw[arr] (ethonly) -- (acc);
  \draw[arr] (start) -- node[lbl, right]{port 255} (po);
  \draw[arr] (start.west) -- (-1.6,0) |- node[lbl, pos=0.25, left]{other} (eth.west);
  \draw[arr] (po) -- (eth);
  \draw[arr] (eth) -- node[lbl]{0x0800} (v4); \draw[arr] (eth) -- node[lbl]{0x86DD} (v6);
  \draw[arr] (eth) -- node[lbl, above]{0x8100 / 0x88A8} (vl); \draw[arr] (vl) |- node[lbl, pos=0.3, left]{IPv4} (v4);
  \draw[arr] (v4) -- node[lbl, left]{6} (tcp); \draw[arr] (v6) -- node[lbl, right]{17} (udp);
  \draw[arr] (v4.south east) -- node[lbl, pos=0.75]{17} (udp.north west); \draw[arr] (v6.south west) -- node[lbl, pos=0.75]{6} (tcp.north east);
  \draw[arr] (tcp) -- (opt); \draw[arr] (opt) -- (tp);
  \draw[arr] (udp) -- (d4); \draw[arr] (udp) -- (d6);
  \draw[arr] (tp) -| node[lbl, pos=0.3, above]{kind 3 (certificate)} (x);
  \draw[arr] (tp) -- node[lbl, right]{1--8 bytes} (short);
  \draw[arr] (tp) -| node[lbl, pos=0.3, above]{$\ge$ 9 bytes} (peek);
  \draw[arr] (short) -- node[lbl, left]{0x16 or a\\ method letter} (inc);
  \draw[arr] (peek) -- node[lbl, right]{0x16 0x03, type 1 / 2} (tlschk);
  \draw[arr] (peek) -| node[lbl, pos=0.3, above]{else} (httpchk);
\end{tikzpicture}
\caption{The parser's top level (\code{parser.p4} with \code{parser\_tls.p4}, \code{parser\_http.p4},
\code{parser\_dhcp.p4}, \code{parser\_x509.p4} included inside it). IPv4 fragments other than the first are not
read beyond IP. Anything not matched goes to \code{accept}: the packet is still forwarded.}
\end{figure}

\begin{figure}[H]\centering
\begin{tikzpicture}
  \node[st] (next) at (0,0) {opt\_next\\ left = 0?};
  \node[st] (kind) at (3.2,0) {opt\_kind: lookahead\\ kind, extract 1 byte,\\ append kind};
  \node[st] (lc) at (6.6,0) {opt\_len\_check:\\ room for the\\ length byte?};
  \node[st] (len) at (9.8,0) {opt\_len: lookahead\\ length, extract};
  \node[st] (pad) at (13.6,0) {opt\_padding:\\ malformed, rest\\ kept, not read};
  \node[st, fill=white] (pay) at (0,-2.4) {tcp\_payload};
  \node[st] (skip) at (3.3,-2.4) {opt\_skip: extract\\ the body byte by byte};
  \node[st] (ws) at (6.5,-2.4) {opt\_wscale: lookahead\\ 8 bits, extract 1};
  \node[st] (mss) at (9.7,-2.4) {opt\_mss: lookahead\\ 16 bits, extract 2};
  \draw[arr] (next) -- node[lbl, above=1mm]{no} (kind);
  \draw[arr] (next) -- node[lbl, right]{yes} (pay);
  \draw[arr] (kind) -- node[lbl, above=1mm]{kind $\ge$ 2} (lc);
  \draw[arr] (kind.north) -- ++(0,0.55) -| node[lbl, pos=0.25]{0 (EOL), 1 (NOP): one byte, continue} (next.north);
  \draw[arr] (lc) -- node[lbl, above=1mm]{yes} (len);
  \draw[arr] (len) -- node[lbl, above=2mm]{len $<$ 2 or\\ too long} (pad);
  \draw[arr] (len) -- node[lbl, right]{2, len 4} (mss);
  \draw[arr] (len) -- node[lbl, pos=0.45]{3, len 3} (ws);
  \draw[arr] (len) -- node[lbl, pos=0.3]{other} (skip);
  \draw[arr] (pad.south) -- (13.6,-3.7) -| (pay.south);
  \draw[arr] (mss.south) -- (9.7,-3.3) -- (-1.7,-3.3) |- (next.west);
  \draw (ws.south) -- (6.5,-3.3); \draw (skip.south) -- (3.3,-3.3);
\end{tikzpicture}
\caption{TCP option walk (\code{parser.p4}): \code{left} = data offset $\times$ 4 $-$ 20 counts the option bytes
still to read, so the walk never reads payload. Every byte is \emph{extracted} (into \code{hdr.tcp\_opt}), so the
deparser re-emits the options unchanged and this walk needs no second pass. EOL is a one-byte option and the walk
continues: FoxIO (through tshark) lists every padding byte after EOL as another 0.}
\label{fig:opt}
\end{figure}

\begin{figure}[H]\centering
\begin{tikzpicture}
  \node[st] (peek) at (0,0) {tls\_peek: 9 bytes:\\ record + handshake header};
  \node[st] (chk) at (4.0,0) {tls\_check\\ hsLen + 9 $\le$ payLen and\\ hsLen + 4 $\le$ recLen};
  \node[st] (hello) at (8.2,0) {tls\_hello: version,\\ random, session id};
  \node[st] (inc) at (4.0,-1.8) {tls\_incomplete};
  \node[st] (ch) at (7.0,-1.8) {ch\_ciphers: value by value,\\ GREASE? capacity?};
  \node[st] (sh) at (10.6,-1.8) {sh\_body: the\\ chosen cipher};
  \node[st] (ext) at (8.2,-3.5) {ext\_block: type, length\\ (appended in order)};
  \node[st] (sk) at (12.2,-3.5) {ext\_skip: the rest\\ of a body};
  \node[st] (sni) at (1.2,-5.3) {SNI (0): present};
  \node[st] (alpn) at (4.6,-5.3) {ALPN (16): first name's\\ length, first, last byte};
  \node[st] (sv) at (8.6,-5.3) {supported\_versions (43):\\ highest non-GREASE};
  \node[st] (sig) at (12.4,-5.3) {signature\_algorithms (13):\\ value by value};
  \draw[arr] (peek) -- (chk); \draw[arr] (chk) -- node[lbl]{yes} (hello);
  \draw[arr] (chk) -- node[lbl, right]{no} (inc);
  \draw[arr] (hello) -- node[lbl]{1} (ch); \draw[arr] (hello) -- node[lbl]{2} (sh);
  \draw[arr] (ch) -- (ext); \draw[arr] (sh) -- (ext);
  \draw[arr] (ext) -- (sni); \draw[arr] (ext) -- (alpn); \draw[arr] (ext) -- (sv); \draw[arr] (ext) -- (sig);
  \draw[arr] (ext) -- (sk);
\end{tikzpicture}
\caption{TLS hello walk (\code{parser\_tls.p4}). Each list value is read with \code{lookahead} and skipped with
\code{advance}; a GREASE value (\code{0x?a?a}, equal bytes) goes back to the loop without being kept; a value
beyond the capacity sets an overflow bit. Every extension type is appended in order (client: without GREASE).}
\end{figure}

\begin{figure}[H]\centering
\begin{tikzpicture}
  \node[st] (m) at (0,0) {method: bytes up to\\ the space};
  \node[st] (t) at (3.6,0) {target: skipped up to\\ the space};
  \node[st] (v) at (7.2,0) {version: the last 8 bytes\\ of the line};
  \node[st] (n) at (7.2,-1.9) {header name: bytes\\ appended to names;\\ lower-cased copy kept};
  \node[st] (end) at (11.2,-1.9) {empty line (CR LF\\ or bare LF):\\ parsed = 1};
  \node[st] (k) at (3.6,-1.9) {is it cookie, referer\\ or accept-language?};
  \node[st] (cv) at (0,-3.8) {Cookie value: pairs\\ split at "; " (a zero\\ byte between pairs)};
  \node[st] (lv) at (3.6,-3.8) {Accept-Language: up to\\ 4 characters before ',',\\ '-' dropped};
  \node[st] (val) at (7.2,-3.8) {any other value:\\ skipped to the line end};
  \draw[arr] (m) -- (t); \draw[arr] (t) -- (v); \draw[arr] (v) -- (n); \draw[arr] (n) -- node[lbl]{':'} (k);
  \draw[arr] (k) -- node[lbl]{cookie} (cv); \draw[arr] (k) -- node[lbl, right]{language} (lv);
  \draw[arr] (k) -- node[lbl]{other} (val);
  \draw[arr] (n) -- node[lbl]{empty} (end);
  \draw[arr] (val.east) -| ($(n.south)+(4mm,0)$);
\end{tikzpicture}
\caption{HTTP request walk (\code{parser\_http.p4}), one byte per state. A name is appended to the names list as
it is read; when it turns out to be Cookie or Referer (excluded from JA4H's name list), the last 6 or 7 bytes are
removed with a constant-width slice. If the payload ends before the empty line, \code{incomplete} is set instead.}
\end{figure}

\begin{figure}[H]\centering
\begin{tikzpicture}[node distance=4mm and 5mm]
  \node[st] (tlv) {\textbf{tlv}: read tag,\\ length (short /\\ 0x8n long form)};
  \node[st, right=14mm of tlv] (ret) {select on\\ \code{meta.x509.ret}};
  \node[st, right=14mm of ret, yshift=14mm] (r1) {R\_CERT, R\_TBS:\\ enter SEQUENCE};
  \node[st, below=of r1] (r2) {R\_VERSION, R\_SERIAL,\\ R\_SIGALG, R\_VALIDITY,\\ R\_SPKI: skip};
  \node[st, below=of r2] (r3) {R\_NAME, R\_SET, R\_ATV:\\ countdowns nameLeft,\\ setLeft, itemLeft};
  \node[st, below=of r3] (r4) {R\_OID: copy\\ $\langle$len$\rangle\langle$bytes$\rangle$ to\\ issuer / subject / exts};
  \node[st, below=of r4] (r5) {R\_EXTSEQ, R\_EXT:\\ extension OIDs};
  \draw[arr] (tlv) -- (ret);
  \foreach \n in {r1,r2,r3,r4,r5} \draw[arr] (ret.east) -- (\n.west);
  \draw[arr] (r3.south west) to[bend left=15] node[lbl]{next element:\\ set \code{ret}, go to tlv} (tlv.south);
\end{tikzpicture}
\caption{DER walk (\code{parser\_x509.p4}). A P4 parser has no calls, so one shared \code{tlv} reader is entered
with a ``return code'' and a \code{select} on it continues in the right state. Nesting ends are known from byte
countdowns per level, so no stack is needed.}
\label{fig:x509}
\end{figure}

%---------------------------------------------------------------------
\subsection{Ingress: the order of the controls}

\begin{figure}[H]\centering
\begin{tikzpicture}
  \node[swbox, text width=52mm] (fw) at (0,0) {inPort := ingress\_port (not on a resubmit);\\ from the controller? drop : \code{forward.apply()}};
  \node[dec] (rs) at (0,-1.5) {second pass?};
  \node[term] (nothing) at (5.4,-1.5) {nothing more: the original\\ is forwarded (and copied)};
  \node[dec] (tcp) at (0,-2.9) {TCP?};
  \node[swbox, text width=40mm] (udp) at (4.6,-2.9) {UDP: \textbf{Dhcp}; walked? resubmit};
  \node[swbox, text width=52mm] (ck) at (0,-4.1) {\textbf{ConnKey}: slot (CRC16), check (CRC32)};
  \node[swbox, text width=52mm] (cid) at (0,-5.0) {\textbf{ConnId}: the connection's number};
  \node[swbox, text width=52mm] (hs) at (0,-5.9) {\textbf{Handshake} (real packets only): JA4T / JA4TS / JA4L};
  \node[dec] (k) at (0,-7.2) {kind?};
  \node[swbox, text width=30mm] (cert) at (-2.7,-8.6) {\textbf{Stream} + \textbf{X509}};
  \node[swbox, text width=30mm] (msg) at (2.4,-8.6) {\textbf{Stream} + \textbf{Tls} + \textbf{Http}};
  \node[swbox, text width=52mm] (ssh) at (0,-9.8) {\textbf{Ssh} (real packets and the flush): counts};
  \node[swbox, text width=52mm] (re) at (0,-10.7) {walked? \code{resubmit\_preserving\_field\_list(0)}};
  \node[swbox, text width=52mm] (hd) at (0,-11.9) {the handshake digest, from one call; postponed\\ to the second pass if this pass sent another digest};
  \draw[arr] (fw) -- (rs); \draw[arr] (rs) -- node[lbl, above]{yes} (nothing);
  \draw[arr] (rs) -- node[lbl, right]{no} (tcp); \draw[arr] (tcp) -- node[lbl, above]{UDP} (udp);
  \draw[arr] (tcp) -- node[lbl, right]{yes} (ck); \draw[arr] (ck) -- (cid); \draw[arr] (cid) -- (hs); \draw[arr] (hs) -- (k);
  \draw[arr] (k) -| node[lbl, above, pos=0.3]{3 (certificate)} (cert); \draw[arr] (k) -| node[lbl, above, pos=0.3]{0 or 1} (msg);
  \draw[arr] (k.east) -- node[lbl, above, pos=0.7]{2 (flush)} (5.2,-7.2) |- ($(ssh.east)+(0,0.1)$);
  \draw[arr] (cert) |- (ssh.west); \draw[arr] (msg) |- ($(ssh.east)+(0,-0.1)$);
  \draw[arr] (ssh) -- (re); \draw[arr] (re) -- (hd);
  \draw[arr] (nothing.east) -- (8.2,-1.5) |- ($(hd.east)+(0,0.12)$);
  \draw[arr] (udp.south) -- (4.6,-11.78) -- ($(hd.east)+(0,-0.12)$);
\end{tikzpicture}
\caption{\code{MyIngress} in \code{main.p4}. ``kind'' is the packet-out kind of a packet from the controller (0
for a real packet): 1 a joined message, 2 an SSH flush (Ssh only), 3 a certificate. Ssh runs after Stream, which
sets the packet's direction (client or server) that Ssh needs on ports other than 22.}
\end{figure}

\why{Every control is applied exactly once and in this order because (a) BMv2 allows a table / control instance
to be applied in one place only, (b) Stream must see the parse result (incomplete, parsed) before Tls/Http decide
to report, and (c) BMv2 keeps \emph{one digest per packet}: the handshake digest is therefore taken out of
Handshake and sent at the end, from one call, and postponed to the second pass when this pass sends a message
digest (the handshake's ACK can carry a whole ClientHello) or Ssh's digest (on the SYN of a reused port, the previous
connection's last window; on an SSH banner piggybacked on the handshake's ACK).}

%---------------------------------------------------------------------
\subsection{Per-connection state}

\begin{figure}[H]\centering
\begin{tikzpicture}[node distance=4mm and 6mm, every node/.style={font=\scriptsize}]
  \node[swbox, text width=40mm] (k) {endpoints in a fixed order:\\ smaller (address, port) first};
  \node[swbox, right=of k, text width=32mm] (h1) {slot = CRC16(ver, a1, p1, a2, p2)\\ mod 65,536};
  \node[swbox, below=of h1, text width=32mm] (h2) {check = CRC32(same)};
  \node[swbox, right=of h1, text width=50mm, yshift=-6mm, minimum height=26mm, align=left] (regs) {\textbf{registers[slot]}\\
    ConnId: check, number, SYN sequence number\\
    Handshake: number, state, tSYN, tSYNACK, TTLs\\
    Stream: number, client side, ISNs, started, pending,\\ \quad pendingSeq, edge, gap, gap open, HTTP seen\\
    Tls: number, done (bit 0 CH, bit 1 SH)\\
    Http: number, seen, last request seq\\
    Ssh: number, 4 counts, batch (8$\times$16 bits), count,\\ \quad banner seen};
  \draw[arr] (k) -- (h1); \draw[arr] (k) |- (h2);
  \draw[arr] (h1) -- (h1 -| regs.west); \draw[arr] (h2) -- (h2 -| regs.west);
  \node[dec, below=12mm of regs] (m) {stored number =\\ connection's number?};
  \node[term, left=10mm of m] (same) {same connection:\\ use the state};
  \node[term, right=6mm of m] (new) {another one: start\\ fresh, take the slot};
  \draw[arr] (regs) -- (m); \draw[arr] (m) -- node[lbl, above]{yes} (same); \draw[arr] (m) -- node[lbl, above]{no} (new);
\end{tikzpicture}
\caption{ConnKey and ConnId (\code{main.p4}) and the registers. Both directions of a connection find the same slot.
ConnId gives the connection a number: a new one when another connection (a different check) comes to the slot, and
when a SYN with a new initial sequence number arrives on the same endpoints (port reuse); a repeated SYN keeps it.
Every control keeps its state with this number, and every digest and copy carries it.}
\end{figure}

\why{A P4 switch cannot insert table entries itself, so per-connection state lives in register arrays indexed by a
hash. CRC16 gives the index; the independent CRC32 tells two connections in one slot apart with a 1 in $2^{32}$
chance of a false match.}
\why{\emph{The connection number}: the endpoints alone do not identify a connection. A client, and much faster a NAT,
reuses addresses and ports for a new connection after the old one closed. Before the number, the second
connection's fingerprints were dropped by the controller (it kept the first value per endpoints) and SSH counts of
the two connections were added together into a wrong JA4SSH (found by \code{test\_port\_reuse}).}
\whynot{\emph{The controller installs one table entry per connection:} each new connection would wait for a
round trip to the controller before its SYN state exists, and the SYN's timestamp (JA4L) would be lost.
\emph{A larger hash:} 65,536 slots is plenty for one link of a test network; the size is one constant.}
